%!TEX root = ../../main.tex
% ============================================================
% 力学作图 / 浮力示意图
% 本文件由 main.tex 通过 \input 引入，请勿单独编译
% 重构日期: 2026-04-11
% ============================================================

\subsection{解答：画出漂浮在液面上的木块所受力的示意图}

\vspace{0.6cm}

\begin{center}
\begin{tikzpicture}[>=Stealth, line join=round, line cap=round, scale=1.2]
  % 1. 绘制烧杯内的液体背景 (灰色填充)
  % 水面在 y=1.5，杯底在 y=0，边缘宽度为 1.8 到 -1.8，底部圆角 R=0.6 更自然
  \fill[gray!40] (-1.8, 1.5) -- (-1.8, 0.6) arc(180:270:0.6) -- (1.2, 0) arc(-90:0:0.6) -- (1.8, 1.5) -- cycle;
  
  % 2. 绘制液体表面线 (从杯壁到杯壁，中间部分会被木块遮挡)
  \draw[black, line width=0.8pt] (-1.8, 1.5) -- (1.8, 1.5);
  
  % 3. 绘制盛装液体的烧杯 (加粗黑线，自然平滑外翻，大圆弧过渡底角)
  \draw[black, line width=1.5pt] (-2.0, 2.5) to[out=-80, in=90] (-1.8, 1.5) -- (-1.8, 0.6) arc(180:270:0.6) -- (1.2, 0) arc(-90:0:0.6) -- (1.8, 1.5) to[out=90, in=260] (2.0, 2.5);
  
  % 4. 绘制漂浮的木块 (中间矩形，纯白填充遮挡后方液体线)
  % 木块中心设在 (0, 1.4)，木块大小为 1.2 * 1.0 (x = -0.6 to 0.6, y = 0.9 to 1.9)
  % 浸入液体深度为 1.5 - 0.9 = 0.6，露出液面为 1.9 - 1.5 = 0.4
  \draw[black, fill=white, line width=1.5pt] (-0.6, 0.9) rectangle (0.6, 1.9);

  % ===================================
  % 一、绘制题目所求的两个力（重力与浮力，等大反向）
  % ===================================
  
  % 标注重心
  \coordinate (Center) at (0, 1.4);
  \fill[black] (Center) circle (2pt);
  
  % 重力向下，加粗黑色矢量线 (长度 1.1)
  \draw[->, black, line width=1.8pt] (Center) -- (0, 0.3) 
    node[right, font=\large\bfseries] {$G$};
    
  % 浮力向上，与重力严格等长，加粗黑色矢量 (长度 1.1)
  \draw[->, black, line width=1.8pt] (Center) -- (0, 2.5) 
    node[right, font=\large\bfseries] {$F_{\text{浮}}$};

\end{tikzpicture}
\end{center}

\vspace{0.6cm}

{\small\color{gray}
\textbf{解题说明：}\\
1. \textbf{受力分析：}木块在水面处于“静止状态”（即漂浮），这是典型的二力平衡情形。因此木块整体上仅受到两个力的作用：一个是地球施加的竖直向下的重力 $G$，另一个是液体对其施加的竖直向上的浮力 $F_{\text{浮}}$。\\
2. \textbf{长度规格（画法关键）：}根据二力平衡约束条件，重力大小必须严格等于浮力大小（即 $|G| = |F_{\text{浮}}|$）。所以在实际作图阅卷中，这两条有向线段的端点长度\textbf{必须用直尺画得严格一样长}。\\
3. \textbf{作用点处理：}重力和浮力的作用点虽然物块受力实际微观部位不同（重力在重心，浮力在浸润体浮心），但为了作图教学方便、清晰表意宏观整体力学状态，通常将两个力共点且等效描点在物体的宏观几何重心位置上。并朝正上和正下方相互拉开对立力矩即可。
}

\vspace{0.6cm}

{\small\color{blue!70!black}
\textbf{绘图基本原理与步骤：}\\
\textbf{1. 容器渲染与层叠：}为还原最纯正的教辅原图样板，优先采取底部涂切逻辑 \texttt{\textbackslash fill[gray!40]} 渲染内壁包围的纯灰色区域作为液体区；直接横穿一根 $y=1.5$ 的水面实线；接着用 \texttt{\textbackslash draw[black]} 运用 \texttt{arc} 圆弧以及出入角向量 \texttt{out=-60, in=90} 等贝塞尔方法精密雕刻出极具自然张力的烧杯容器外侧翻边和杯筒。\\
\textbf{2. 绝妙的视差遮挡属性（Z-Index 运用）：}此处的木质矩形 \texttt{rectangle} 因为设定上就悬跨在浮潜截面线上，因此只需顺势为其添加上 \texttt{fill=white} 强力纯白填充遮罩属性，并将绘制层次（代码行数）在容器水体代码的下文处执行。系统渲染渲染器便会自动将该区块上方覆盖的基准水线、背景全副物理阻挡遮盖！轻松复活物理实体的三维感，省去庞大的手工切线代码。\\
\textbf{3. 二力等大的数学绑定：}重点在矢量表现端处，采用平行的固定距离数学位移坐标实现。源点坐标设在极简的原系对称中心 \texttt{(0, 1.4)} 上，那么拉升到 $y=2.5$ （长度 1.1）的力必与拉伸至下位 $y=0.3$ 的重力严格具有 1.1 的位移值。规避一切视觉透视歪斜或人眼误差，彻底合乎中学标量静定要求。
}

\vspace{1.0cm}

%% ==============================
%% 第7题：画出密度计在酒精中的大致位置
%% ==============================

\newpage

\subsection{解答：画出密度计在酒精中的大致位置}

\vspace{0.6cm}

\begin{center}
\begin{tikzpicture}[>=Stealth, line join=round, line cap=round, scale=1.2]
  % ---------------------------------------------------
  % (a) 水中情形
  % ---------------------------------------------------
  \begin{scope}[shift={(-2.2, 0)}]
    % 1. 液体背景 (水，灰色填充)
    % 水面在 y=2.0，杯底在 y=0 
    \fill[gray!40] (-1.2, 2.0) -- (-1.2, 0.4) arc(180:270:0.4) -- (0.8, 0) arc(-90:0:0.4) -- (1.2, 2.0) -- cycle;
    
    % 2. 液体表面线 
    \draw[black, line width=0.8pt] (-1.2, 2.0) -- (1.2, 2.0);
    
    % 3. 烧杯轮廓
    \draw[black, line width=1.5pt] (-1.4, 2.5) to[out=-80, in=90] (-1.2, 1.8) -- (-1.2, 0.4) arc(180:270:0.4) -- (0.8, 0) arc(-90:0:0.4) -- (1.2, 1.8) to[out=90, in=260] (1.4, 2.5);
    
    % 4. 密度计 (在水中)
    % 假设水深 2.0，设浸入深度 h1 = 1.2，则底端在 y=0.8
    % 本体使用纯白填充以遮挡背后的水位实线
    \draw[black, fill=white, line width=1.2pt] (-0.2, 0.8) rectangle (0.2, 4.3);
    
    % 密度计内部配重物（底部铅粒/铁丝示意）
    \fill[gray!60] (-0.15, 0.85) rectangle (0.15, 1.1);
    \draw[black, line width=0.6pt] (-0.08, 1.1) to[out=60, in=-120] (0.08, 1.4) to[out=60, in=-120] (-0.08, 1.7) to[out=60, in=-120] (0.0, 2.0);
    
    % 5. 标注
    \draw[black, line width=0.8pt] (1.3, 0.8) -- (2.0, 0.8) node[right, font=\large] {水};
    \node at (0, -0.8) {\Large (a)};
  \end{scope}

  % ---------------------------------------------------
  % (b) 酒精中情形 (作图答案)
  % ---------------------------------------------------
  \begin{scope}[shift={(2.2, 0)}]
    % 1. 液体背景 (酒精，灰色填充)
    \fill[gray!40] (-1.2, 2.0) -- (-1.2, 0.4) arc(180:270:0.4) -- (0.8, 0) arc(-90:0:0.4) -- (1.2, 2.0) -- cycle;
    
    % 2. 液体表面线 
    \draw[black, line width=0.8pt] (-1.2, 2.0) -- (1.2, 2.0);
    
    % 3. 烧杯轮廓
    \draw[black, line width=1.5pt] (-1.4, 2.5) to[out=-80, in=90] (-1.2, 1.8) -- (-1.2, 0.4) arc(180:270:0.4) -- (0.8, 0) arc(-90:0:0.4) -- (1.2, 1.8) to[out=90, in=260] (1.4, 2.5);
    
    % 4. 密度计 (在酒精中，使用稍粗黑线作为本题学生需绘制的作图答案)
    % 设酒精的密度约为水的 0.8 倍，由于整体要求浮力一致（均等于重力），排开液体体积必须成反比增大。
    % 所以在酒精中浸入深度 h2 = h1 * (1.0/0.8) = 1.25 * 1.2 = 1.5。即浸入深度加深 0.3。
    % 则密度计底端坐标应该设定为 y = 2.0 - 1.5 = 0.5。对应顶端下降为 y = 4.0。
    \draw[black, fill=white, line width=1.8pt] (-0.2, 0.5) rectangle (0.2, 4.0);
    
    % 密度计内部配重物同步下移
    \fill[gray!60] (-0.15, 0.55) rectangle (0.15, 0.8);
    \draw[black, line width=0.6pt] (-0.08, 0.8) to[out=60, in=-120] (0.08, 1.1) to[out=60, in=-120] (-0.08, 1.4) to[out=60, in=-120] (0.0, 1.7);
    
    % 5. 标注
    \draw[black, line width=0.8pt] (1.3, 0.8) -- (2.0, 0.8) node[right, font=\large] {酒精};
    \node at (0, -0.8) {\Large (b)};
  \end{scope}

\end{tikzpicture}
\end{center}

\vspace{0.6cm}

{\small\color{gray}
\textbf{解题说明：}\\
1. \textbf{浮沉条件与浮力常性：}简易密度计无论是在水中还是在酒精中，都能够处于稳定的“漂浮状态”。根据力学平衡的静力学方程可以得知：漂浮情况下其受到的竖直向上的浮力 $F_{\text{浮}}$ 必须等于其自身的重力 $G$。即 $F_{\text{浮水}} = F_{\text{浮酒}} = G$。\\
2. \textbf{排开体积与液体密度的反比例解析：}根据阿基米德浮力公式 $F_{\text{浮}} = \rho_{\text{液}} g V_{\text{排}} = \rho_{\text{液}} g S h_{\text{浸}}$。因为纯酒精的密度小于普通水的密度（$\rho_{\text{酒精}} < \rho_{\text{水}}$），在保证产生的浮力大小相同、自身横截面积 $S$ 全局固定的前提下，密度计在酒精中排开这种液体的体积刻度必须强行增大补偿，即在酒精中浸没的深度 $h_{\text{浸}}$ 必定比在水中更深。\\
3. \textbf{作图结论：}因此在绘制图 (b) 酒精中的密度计时，不仅形状要保持完全一致，最关键的考察点在于它的上下坐标应当集体“下沉”。其浸入酒精液面下的细长部分，要远深于浸入水中的长度。
}

\vspace{0.6cm}

{\small\color{blue!70!black}
\textbf{绘图基本原理与步骤：}\\
\textbf{1. 环境等比例封装与镜像复用：}由于 (a) 图和 (b) 图属于同一规格下截然不同的物理容器态，利用 TikZ 内置的局部 \texttt{scope} 嵌套配合 \texttt{shift=\{(-2.2, 0)\}} 以及 \texttt{shift=\{(2.2, 0)\}} 横移逻辑，将左右图分块平移组合。彻底避免了多容器情况下的复杂坐标重构。\\
\textbf{2. 密度计管身的遮挡图层处理：}如上题般如出一辙地，密度计的身管主体 \texttt{rectangle} 参数中加入了 \texttt{fill=white}。这简单粗暴却极具效果地将后方横移截断穿梭的水面波浪实线遮掩，严格契合水面与漂浮物体之间的真实“三维相干与物理透视阻挡”。另外由于密度计量器底部通常缠绕有细铁丝或者配重小铅球以保证直立，我运用 \texttt{rectangle} 和多重交替的贝塞尔方程模拟了其细铁丝纹路（水和酒精中作了随动的相对下移复原），大大增强了还原考试配图的教辅正规感。\\
\textbf{3. 高精度的浮力下沉参数化：}在定量上并不通过任意瞎比划来体现“下沉”：先验地设定了水深 $H=2.0$ 相对参照，刻画其在水内浸没 $h_1 = 1.2$。严格依据高中物理定则取 $\rho_{\text{酒}} \approx 0.8\text{ g/cm}^3$，由此确凿算出如果等浮力该计身在酒精中必会下沉至深度 $h_2 = h_1 \times (1.0 \div 0.8) = 1.5$ 处。从而代码直接敲定了在酒精中该柱体底部的局部标高必然精准停留在 $y = 2.0 - 1.5 = 0.5$；这套严谨的数字逻辑完全避免了凭空人眼主观定位失察的问题，真正用代码逻辑“算出了结果图”而非仅为图做图！
}

\vspace{1.0cm}

%% ==============================
%% 第8题：画出浸没在水中的小球所受浮力的示意图
%% ==============================

\newpage

\subsection{解答：画出浸没在水中的小球所受浮力的示意图}

\vspace{0.6cm}

\begin{center}
\begin{tikzpicture}[>=Stealth, line join=round, line cap=round, scale=1.2]
  % 1. 绘制水面 (实线)
  \draw[black, line width=1.0pt] (-2.5, 2.5) -- (2.5, 2.5);

  % 2. 绘制水的虚线背景 (表示水)
  \begin{scope}
    % 利用 clip 限制虚线不出界
    \clip (-2.5, 0) rectangle (2.5, 2.5);
    % 奇数行虚线
    \foreach \y in {0.2, 0.6, ..., 2.4} {
      \draw[black, dashed, dash pattern=on 4pt off 4pt, line width=0.4pt] (-2.5, \y) -- (2.5, \y);
    }
    % 偶数行错开相位虚线，这更符合典型的流体题干画风
    \foreach \y in {0.4, 0.8, ..., 2.4} {
      \draw[black, dashed, dash pattern=on 4pt off 4pt, dash phase=4pt, line width=0.4pt] (-2.5, \y) -- (2.5, \y);
    }
  \end{scope}

  % 3. 绘制完全浸没的小球 (实心白底黑边圆形)
  % 为了不让水的虚线穿透小球，这里为小球填充了不透明纯白色（fill=white）
  \coordinate (BallCenter) at (0, 1.2);
  \draw[black, fill=white, line width=1.2pt] (BallCenter) circle (0.5);

  % ===================================
  % 一、 绘制题目所求的答案力（浮力）
  % ===================================
  
  % 标注重心
  \fill[black] (BallCenter) circle (2pt);
  
  % 浸没在水中的物体受到的浮力方向竖直向上，作用点在重力几何中心
  \draw[->, black, line width=1.8pt] (BallCenter) -- (0, 3.2) 
    node[right, font=\large\bfseries] {$F_{\text{浮}}$};

\end{tikzpicture}
\end{center}

\vspace{0.6cm}

{\small\color{gray}
\textbf{解题说明：}\\
1. \textbf{确认受力物体与浮力三要素：}题目要求画出“小球”所受“浮力”的示意图。浮力产生的原因是液体对物体上下表面的压力差，根据阿基米德原理及其受力特性，物理学中规定浮力的\textbf{方向始终是竖直向上}的。\\
2. \textbf{确定作用点：}浮力的等效作用点称为“浮心”，对于材质均匀且完全浸没的规则形体（如本题中的小球），浮心与其几何中心（重心）重合。因此，我们将画笔的起点（受力点）点在小球的圆心处。\\
3. \textbf{作图规范：}从圆心开始，沿着竖直向上的方向画一条带实心箭头的线段，表示浮力向量，并在箭头端部附近标注大写字母符号 $F_{\text{浮}}$ 即可完成规范解答。
}

\vspace{0.6cm}

{\small\color{blue!70!black}
\textbf{绘图基本原理与步骤：}\\
\textbf{1. 水域环境流体构建：}基于顶部的 \texttt{draw} 绘制水面实线界限；随之通过 TikZ 的循环语法 \texttt{\textbackslash foreach} 分奇数和偶数阵列配合 \texttt{dashed} 和错位配置 \texttt{dash phase=4pt} 参数，快速铺展出等间距且交错的精美点阵横向虚线群，高度还原了物理原题教辅中典型的纯正标量“水池”虚线质感。\\
\textbf{2. 实体遮挡技巧：}利用 \texttt{\textbackslash draw[fill=white]} 语法绘制小球。这不但能够勾勒出球体主边界 \texttt{circle(0.5)}，更关键的是白底不透明度能在渲染层级上自然“强势掩盖遮挡”下方的水纹穿梭虚线群，从而彻底避免了水层虚假穿透球体内部的低级失真视觉错位状况。\\
\textbf{3. 力学标定：}单独抽出统一的作用点几何坐标 \texttt{coordinate (BallCenter)} 声明。直接在球心用 \texttt{fill} 打上实心黑点标记重心，并在视觉的顶层引出加宽加粗并附带 \texttt{Stealth} 实心箭头的竖直向上长力矩射线，契合标准作图题答案的高精读样张设计。
}

\vspace{1.0cm}

%% ==============================
%% 第9题：画出试管转到竖直位置时的液面
%% ==============================

\newpage

\subsection{解答：画出试管转到竖直位置时的液面}

\vspace{0.6cm}

\begin{center}
\begin{tikzpicture}[>=Stealth, line join=round, line cap=round, scale=1.2]
  
  % ===================================
  % 一、 绘制原题干几何结构（黑色实线与虚线）
  % ===================================

  % 1. 绘制水平地面及阴影
  \draw[black, line width=1.0pt] (-2.0, 0) -- (4.0, 0);
  \foreach \x in {-1.8, -1.4, ..., 3.8} {
    \draw[black, line width=0.6pt] (\x, 0) -- (\x-0.2, -0.2);
  }

  % 2. 绘制位置 1 的试管（倾斜）和内部初始水
  \begin{scope}[shift={(0.35, 0.25)}, rotate=-40]
    % 填充位置 1 的水 (保持水平)
    % 水柱中心在 L=2.0 处。全局水平面在局部表现为斜率为 tan(40)
    \fill[gray!50] (-0.25, {2.0 - 0.25*tan(40)}) -- (-0.25, 0) arc(-180:0:0.25) -- (0.25, {2.0 + 0.25*tan(40)}) -- cycle;
    % 水面轮廓线
    \draw[black, line width=1.0pt] (-0.25, {2.0 - 0.25*tan(40)}) -- (0.25, {2.0 + 0.25*tan(40)}); 
    
    % 绘制位置 1 的试管外壳 (实线)
    \draw[black, line width=1.2pt] (-0.25, 4.0) -- (-0.25, 0) arc(-180:0:0.25) -- (0.25, 4.0) -- cycle;
  \end{scope}

  % 标注"1"：处于试管1右侧水面附近
  \node at (2.6, 2) {\Large 1};

  % 3. 绘制位置 2 的试管（竖直，虚线）
  \begin{scope}[shift={(0, 0.25)}]
    % ===================================
    % 二、 绘制题目所求的答案液面部分
    % ===================================
    % 倾斜时水的中心线在 L=2.0，竖直时其高度仍为 L=2.0。此时液面高度必定高于倾斜时。
    \fill[gray!50] (-0.25, 2.0) -- (-0.25, 0) arc(-180:0:0.25) -- (0.25, 2.0) -- cycle;
    
    % 加粗强调竖直时的液面（本题作图答案），并保持原题风格，使用明显稍粗的黑线
    \draw[black, line width=2.0pt] (-0.25, 2.0) -- (0.25, 2.0);
    
    % 绘制试管轮廓 (虚线)
    \draw[dashed, black, line width=1.0pt] (-0.25, 4.0) -- (-0.25, 0) arc(-180:0:0.25) -- (0.25, 4.0) -- cycle;
    
    \node at (0.4, 4.4) {\Large 2};
  \end{scope}

  % 4. 添加转动标志（弯曲箭头）
  \draw[->, black, line width=1.0pt] (2.0, 3.6) to[out=130, in=20] (0.5, 4.2);

\end{tikzpicture}
\end{center}

\vspace{0.6cm}

{\small\color{gray}
\textbf{解题说明：}\\
1. \textbf{作图分析（画出液面位置）：}试管内水的质量和体积是不变的，设试管的底面积为 $S$。当试管倾斜搁置在位置 1 时，水面是水平的，由于圆柱斜截体体积定理，它等效于将水柱“拉平”后齐平在其几何中心处。设这管水若竖直拉平时的管内长度为 $L$。因此，当试管立起到位置 2 时，液柱沿管壁长度恰好是原来位置 1 液面中心的所在管壁长度（即中心点仍然在内部坐标 $L$ 处）。由于试管位置 2 完全竖直，虽然水在管内所占的“长度刻度”不变，但其中心相对于桌面的\textbf{竖直高度相对明显增高}（原高度等于所在管壁长度乘以倾角对应的正弦或余弦分量），直观表现为：图上该水面直线相较于位置 1 水面的中央具有明显的直观上升。\\
2. \textbf{填空（压强变化）：}填空答案为\textbf{变大}。根据液体压强公式 $p = \rho g h$，其中 $h$ 为液体深度（液面到管底的垂直高度）。在此管由斜起立倒直转动的过程中，其内水体也整体被抬起，最高面（液面）对最低面（试管底）的垂直直线深度 $h$ 明显增大。因此，水对试管底部的压强\textbf{变大}。
}

\vspace{0.6cm}

{\small\color{blue!70!black}
\textbf{绘图基本原理与步骤：}\\
\textbf{1. 结构复刻：}利用 TikZ 的 \texttt{foreach} 遍历坐标绘制等距斜短线阵列，构建带阴影的基准水平面环境。\\
\textbf{2. 旋转坐标系与切面一致性质保障：}为确保倾斜试管无论旋转多大角度，底部的半圆极点始终“悬挂并相切”于桌面，不能简单依赖从管口等处向下方旋转探测位置。需要将局部 \texttt{scope} 的中心原点（旋转锚点）精准设置在距离地面高度正好等于试管圆倒角半径 $R=0.25$ 的位置：\texttt{shift=\{(x\_c, 0.25)\}}。以此点为锚固圆心执行 \texttt{rotate=-40} 的切向倾角自变换后，绘制出底面半圆曲线的物理最低点依然始终严格切附在地面桌上，此思路规避了手工直接的推算复杂切面坐标的盲点。\\
\textbf{3. 水位倾斜的自适应建模与守恒体积体现：}在自身处于倾斜姿态的局部坐标系 \texttt{scope} 内部视图下，原本全局的重力“水平面”变成了有确定斜率的倾斜直线，斜率为全局转角即 pgfmath 解析值 \texttt{tan(40)}。基于流体特性，我们只需固定体积约束下常量管中水长 $L=2.0$ 作为动态水面的纵截面中心点集；直接借助公式表达式 \texttt{y = 2.0 $\pm$ R*tan(40)} 便能精确求取倾斜水面的左右两侧切点相对坐标并执行安全闭合 \texttt{cycle} 进行液域着色；对应于受力终态也就是位置 2 （完全竖直），出于体积等价等效公理，该终态液面高度可以直接直接的切落截断于相对常量 \texttt{y=2.0} 。辅以加重加粗的纯黑黑实线，这种直接渲染能够使该体积逻辑高度无论在参数变化中均稳定精确大于倾斜状态时的物理实际视觉中心标高，真正实现在原汁原味代码化和物理规律的双重自洽检验。
}

%% ==============================
%% 第18题：画出下列各力的示意图
%% ==============================
